\section{Extremal lattice sections and design moments}
\label{sec:sections}

The preceding chapters specify each layer and check its products. Here we
select vectors from a large known lattice shell and obtain lower-dimensional
configurations by projection. The required information is the number of
vectors with particular projection signatures. Dimension 43 counts shell
vectors perpendicular to a five-dimensional section; dimension 45 projects
three fibres of a three-dimensional section. Both become finite multiplicity
problems.

\subsection{Mother shells, projection domains and
moments}

An extremal even unimodular lattice in dimension 48 has minimum squared
norm six. Its theta series is

\[
E_4^6-1440E_4^3\Delta+125280\Delta^2
=1+52416000q^3+O(q^4).
\]

Thus its minimal shell has 52416000 vectors. Here
\(E_4=1+240q+2160q^2+6720q^3+\cdots\) and
\(\Delta=q-24q^2+252q^3+\cdots\); the constant term and vanishing first
two coefficients determine the expression in weight 24
\cite{conway1999sphere}. Venkov's theorem makes this shell a spherical
11-design~\cite[Theorem 4.1]{nebe-designs}.

A spherical 11-design has the same averages as the sphere of equal radius
for every polynomial of degree at most eleven. Thus its low-degree moments
along any chosen direction are known without listing its vectors. For example,
\[
\sum_{x\in X}(x\cdot t)^2=\frac{6|X|}{48}\|t\|^2,\qquad
\sum_{x\in X}(x\cdot t)^4=\frac{3\cdot6^2|X|}{48\cdot50}\|t\|^4.
\]
Grouping vectors by projection turns the left sides into linear combinations
of unknown multiplicities, with completely known right sides.

For section vectors \(t_i\) with Gram matrix \(H\), a shell vector has
signature \(y_i=x\cdot t_i\) and projection norm
\(y^{\mathsf T}H^{-1}y\le6\). The signatures are integral when the
section vectors lie in the lattice. Inner-product bounds against known
short section vectors further restrict this finite domain. For degree
\(2m\le10\),

\[
\sum_{x\in X}\prod_{\ell=1}^{2m}(x\cdot t_{i_\ell})
=\frac{|X|6^m}{48\cdot50\cdots(48+2m-2)}
\sum_{\mathcal P}\prod_{(a,b)\in\mathcal P}H_{i_ai_b},
\]

where \(\mathcal P\) ranges over pairings of the labelled positions;
degree zero gives the shell size. Unknown signature multiplicities
therefore satisfy rational linear equations. In dimension 43 these
equations determine the target statistic exactly. In dimension 45 they
are combined with nonnegativity to bound the statistic from below.

\subsection{Dimension 43: counting a child section through its parent}

Let \(X\) be the minimal shell of \(P_{48n}\) \cite{latticecatalogue}. A supplied integer matrix
embeds \(\sqrt3E_8\) into this lattice. Its first five simple roots span
a \(\sqrt3D_5\) section \(U\), and \(X\cap U^\perp\) lies in dimension
43.

The root numbering is fixed by the following standard coordinates:
\[
\alpha_1=\tfrac12(1,-1,-1,-1,-1,-1,-1,1),\qquad
\alpha_2=e_1+e_2,\qquad
\alpha_j=e_{j-1}-e_{j-2}\quad(3\le j\le8).
\]
The supplied $8\times48$ integer matrix $T$ and ambient Gram matrix $G$
satisfy $(TGT^{\mathsf T})_{ij}=3\alpha_i\cdot\alpha_j$; rows
$1,2,3,4,5$ specify the child. This fixes the embedding as well as the
abstract $D_5$ root type. The use of an embedded $\sqrt3E_8$ parent is
developed in Dorofeev--Sun--Wang~\cite{dorofeev}; here its projection
data are used to count this $D_5$ complement.

\begin{proposition}[The specified $D_5$ complement]
\label{prop:section43}
For the matrices $G,T$ supplied in the certificate, the shell
$X\cap U^\perp$ contains exactly $2553792$ vectors. Consequently
$K(43)\ge2553792$.
\end{proposition}

No direction change is needed. Distinct minimal vectors satisfy
$\|x-y\|^2\ge6$, so $x\cdot y\le3$; division by $\sqrt6$ gives a kissing
configuration. The task is to count $|X\cap U^\perp|$ exactly.

For a minimal vector, its projection onto the parent section is
\(\lambda/\sqrt3\) with \(\lambda\in E_8\). The norm bound and the inner
products with every embedded root give a complete domain of 26,641
possible signatures. Of these, 240 are exceptional signatures of the
section roots themselves.

This domain exhausts the possible projections. Nonexceptional $\lambda$
satisfy $\|\lambda\|^2\le18$ and $|\lambda\cdot r|\le3$ for every $E_8$
root $r$. Enumeration gives shells of squared norm 0, 2, 4, 6 and 8,
of sizes 1, 240, 2160, 6720 and 17280 respectively. Adding the 240
exceptions $\lambda=3r$ gives the complete domain.

The extremal 48-dimensional shell is an 11-design. For projection
coordinates, every even moment through degree ten is therefore the
corresponding spherical moment. Averaging over the parent Weyl group
yields multiplicities

\[
1092000,\quad116235,\quad9180,\quad495,\quad63/4
\]

on its five nonexceptional norm orbits. As orbit averages of integer
fibre sizes, these multiplicities can be rational.

The perpendicular child has \(1,12,6,24,0\) signatures on those orbits
and twelve exceptional roots. Hence the average child count is

\[
1092000+12(116235)+6(9180)+24(495)+12=2553792.
\]

Thus at least one child has at least 2,553,792 points. To count the
specified child exactly, average only over a subgroup preserving that
child. Explicitly, use the group generated by $-I$, the reflections in
$\alpha_1,\ldots,\alpha_5$, and the reflections in the twelve $E_8$
roots perpendicular to all five. The latter roots form $A_3$; each
reflection is $s_r(x)=x-(x\cdot r)r$. These formulas determine the group
without reference to an ambient automorphism. The complete domain has
43 orbits under this group. Rational moment rows
form a matrix \(A\), with right side \(b\) and target \(c\). The
certificate gives \(w\) satisfying

\[
A^tw=c,\qquad w^tb=2553792.
\]

This proves Proposition~\ref{prop:section43}. The subgroup need not act by
automorphisms of the whole mother lattice: it preserves the signature
domain, the moment equations, and the target statistic.

The calculation determines a linear statistic rather than every fibre.
Even if $Av=b$ has several solutions, $c^{\mathsf T}v$ is constant on
them whenever $c^{\mathsf T}$ lies in the row space of $A$.
Independent moment rows can first be selected by modular linear algebra,
followed by a rational solution for the multiplier. Reproduction then
requires just the two displayed rational identities. The complete
signature domain remains essential: the moments on the right describe
the whole shell, while orthogonality belongs in the target vector $c$,
not in a preliminary deletion of all other fibre columns.

The comparison section has the same abstract \(D_5\) Gram type. What
changes is the embedding, which can change the perpendicular population.
Its extension to an \(E_8\) parent supplies additional projection
constraints that certify this population. The additional equations
determine a count of the fixed set; they do not change that set.

\subsection{Dimension 45: 298 witnesses control the projected count}

Use a specified quadratic-residue even neighbour of rank 48. The
code-distance and parity argument in Appendix~\ref{app:ambient} proves its minimum squared norm
six. Its shell again has 52,416,000 vectors and is an 11-design.

Three section vectors have Gram matrix

\[
H=8I_3-2J_3,\qquad H^{-1}=(I_3+J_3)/8.
\]

Let $f(y)$ count shell vectors of signature $y$. The nonexceptional
domain is explicitly
\begin{equation}
\begin{split}
D_0=\{y\in\mathbb Z^3:\ &|y_i|\le3,\quad |y_1+y_2+y_3|\le3,\\
&\textstyle\sum_i y_i^2+(\sum_i y_i)^2\le48\}.
\end{split}
\label{eq:domain45}
\end{equation}
The section roots are $\pm t_1,\pm t_2,\pm t_3$ and
$\pm(t_1+t_2+t_3)$. Inner products with these eight vectors give the
linear bounds in~\eqref{eq:domain45}, while $y^{\mathsf T}H^{-1}y\le6$
gives its quadratic bound. Their own signatures form
$E=\{\pm He_1,\pm He_2,\pm He_3,\pm H(1,1,1)^{\mathsf T}\}$.
Each has multiplicity one: its section projection already has squared
norm six. Enumeration of the displayed three-coordinate domain gives
$|D_0|=231$ and $|D_0\cup E|=239$.

For each pair $\{y,-y\}\subset D_0$, take the lexicographically smaller
representative and set $v_{[y]}=f(y)=f(-y)$; retain zero separately.
There are $(231+1)/2=116$ variables. For a multi-index $a$ of even
total degree at most ten, its moment row is
\[
A_{a,[y]}=(2-\mathbf1_{y=0})\prod_i y_i^{a_i},\qquad
b_a=M_a-\sum_{y\in E}\prod_i y_i^{a_i},
\]
where $M_a$ is the spherical moment defined above. There are
$\sum_{j=0}^5\binom{2j+2}{2}=161$ such monomials. This constructs all
entries of the $161\times116$ system $Av=b$.

The fibre count $f(y)$ counts shell vectors satisfying
$(x\cdot t_1,x\cdot t_2,x\cdot t_3)=y$.
The target is $f(e_1)+f(e_2)+f(e_3)$, not every individual multiplicity.
After moving known exceptional contributions to the right, write the
moment equations as $Av=b$, with $v\ge0$. For a statistic
$c^{\mathsf T}v$, a rational vector $w$ with $c-A^{\mathsf T}w\ge0$ gives
\[
c^{\mathsf T}v=w^{\mathsf T}b+(c-A^{\mathsf T}w)^{\mathsf T}v
\ge w^{\mathsf T}b.
\]
This is the form of the moment certificate. Choosing the statistic well
makes an accessible fibre appear with positive coefficient in the lower
bound, so a small explicit witness set becomes useful.

\begin{proposition}[A positive fibre contribution]
\label{prop:moment45}
The signature multiplicities satisfy

\[
f(e_1)+f(e_2)+f(e_3)
\ge7377408+9f(-2,-2,3).
\]
\end{proposition}
\begin{proof}
Give $c$ coefficient $+1$ on the three classes $[e_i]$, coefficient
$-9$ on $[(-2,-2,3)]$, and zero elsewhere. Order the representatives
and multi-indices lexicographically. The multiplier in
\texttt{dual.json}, supplied in Appendix~\ref{app:data}, then has 161
rational entries and satisfies, by exact matrix multiplication,
$c-A^{\mathsf T}w\ge0$ and $w^{\mathsf T}b=7377408$.
The preceding identity with $v\ge0$ gives the assertion.
\end{proof}

\begin{corollary}[The 45-dimensional construction]
\label{cor:section45}
The specified quadratic-residue neighbour and section give
$K(45)\ge7380090$.
\end{corollary}
\begin{proof}
The certificate supplies 298 distinct integer lattice vectors of squared
norm six and signature \((-2,-2,3)\). Thus the right side is at least
7,380,090.

Project the three fibres \(e_1,e_2,e_3\) onto the section's orthogonal
complement. The projection formula is
$\pi(x)=x-\operatorname{proj}_{U}x$, giving
\[
\pi(x)\cdot\pi(x')=x\cdot x'-y^{\mathsf T}H^{-1}y'.
\]
For $y=e_i$, the squared section projection norm is $1/4$; for distinct
$e_i,e_j$, their section projections have product $1/8$. Since distinct
mother-shell vectors have product at most three, each projected squared
norm is \(23/4\). Products are at
most \(11/4\) within a fibre and \(23/8\) between different fibres.
Dividing by \(\sqrt{23/4}\) gives a kissing configuration:

\[
K(45)\ge7377408+9(298)=7380090.
\]
\end{proof}

A witness search in a Pless lattice gave 292 vectors and the bound
7,380,036. Enlarging the choice of anchors led to 298 vectors in a
quadratic-residue even neighbour, with each additional witness contributing
nine points to the projected bound. To organize the search, fix a
norm-eight anchor $p\in\Lambda$ and set
\[
A_p=\{x\in X:x\cdot p=4\}.
\]
This parent set has exactly 2256 points, with completeness certified by
low-degree moments. Since $x\pm p$ are nonzero lattice vectors,
minimum six gives $|x\cdot p|\le4$. The polynomial
\[
P(t)=t^2(t^2-1)(t^2-4)(t^2-9)
\]
vanishes at all allowed integral products except $\pm4$, where it
equals 20160. The eighth spherical moment gives
$\sum_{x\in X}P(x\cdot p)=90961920$. Antipodal symmetry makes the
two endpoint fibres equal, hence
$|A_p|=90961920/(2\cdot20160)=2256$.
The involution $x\mapsto p-x$ preserves $A_p$. The graph vertices are
its 1,128 pairs \(\{x,p-x\}\) centred at \(p/2\). Two pairs are adjacent
when representatives have inner product one or three; exchanging an
endpoint preserves this condition. Fibre searches become
common-neighbour calculations in this graph.

Use the integer numerator coordinates $p=y/\sqrt{12}$ of the appendix.
The selected $y$ has six coordinates of absolute value three and 42 of
absolute value one, so $y^2=96$. This odd-coordinate shape explores
another part of the neighbour rather than repeating anchors from its
integral sector. The resulting graph has degree 368 and a selected
nonadjacent pair with 70 common neighbours.
The difference \(368-70=298\) identifies the witness
fibre. The correspondence is explicit. Denote the selected endpoints by
$r,s$ and use the section $T=(-s,s-p,p-r)$. For a pair adjacent to the
vertex of $r$ but not that of $s$, choose its endpoint $a$ with
$r\cdot a=3$; then $s\cdot a=2$. Set $w=p-a$. It satisfies
\[
\|w\|^2=6,\qquad (-s\cdot w,(s-p)\cdot w,(p-r)\cdot w)=(-2,-2,3).
\]
Different pairs give different witnesses, so the 298 selected graph
vertices yield precisely the required lattice vectors. The
parent-graph checker in Appendix~\ref{app:data} reconstructs the section and exactly the 298 vectors in the
final certificate. The rational moment inequality then gives the stated
bound.
